%==============================================================================
% فصل چهارم: یافته‌ها
% فقط یافته‌ها — بدون تفسیر تفصیلی (تفسیر در فصل پنجم)
% از درج هم‌زمان یک یافته به‌صورت جدول و نمودار خودداری شود.
%==============================================================================
\ChapterStart
\chapter{یافته‌ها}
\label{ch:results}

\section{مقدمه فصل}
\label{sec:ch4-intro}


فصل پیش رو دارای چند بخش مجزا می باشد که در بخش اول توصیف بیماران مورد مطالعه از نظر ویژگی های جمعیت شناختی، بالینی و میزان پایبندی به اقدامات پیشگیری ثانویه غیر دارویی و در بخش بعد، بررسی ارتباط متغیرهای جمعیت شناختی و بالینی با پایبندی به اقدامات پیشگیری ثانویه غیر دارویی خواهیم پرداخت.


\section{ویژگی‌های جمعیت‌شناختی شرکت‌کنندگان}
\label{sec:demographics}


در این مطالعه 150 بیمار مبتلا به انفارکتوس میوکارد مورد ارزیابی قرار گرفتند. از این تعداد، 105 بیمار (70 درصد) مرد بودند. میانگین $\pm$ انحراف معیار سنی بیماران $62.1 \pm 10.7$سال بود. وضعیت تاهل در این مطالعه در دو دسته متاهل (۱۲۰ نفر، 80 درصد) و مجرد، مطلقه‌ یا بیوه (30 نفر، 20 درصد) بررسی شد. از نظر میزان تحصیلات، ۱۱۹ نفر ($79.3$ درصد) دارای تحصیلات دیپلم و پایین‌تر بودند و ۳۱ نفر ($20.7$ درصد) تحصیلات دانشگاهی داشتند. ارزیابی وضعیت اشتغال نشان داد که 84 نفر از شرکت‌کنندگان (۵۶ درصد) در گروه بیکار، بازنشسته یا خانه‌دار قرار داشته و مابقی شاغل بودند. همچنین از لحاظ محل سکونت، ۸۹ نفر ($59.3$ درصد) در مناطق شهری و مابقی در مناطق روستایی زندگی می‌کردند. در خصوص وضعیت اقتصادی، 65 نفر وضعیت ضعیف ($43.3$ درصد) بود و 85 نفر ($56.6$ درصد) وضعیت اقتصادی متوسط/خوب داشتند. تمامی بیماران مورد مطالعه (۱۵۰ نفر، ۱۰۰ درصد) تحت پوشش بیمه درمانی قرار داشتند.

\section{یافته های بالینی بیماران}


با توجه به توزیع غیر نرمال داده‌های متغیر " تعداد روزهای سپری شده از سکته قلبی" ، از میانه و دامنه میان‌چارکی برای توصیف داده‌ها استفاده شد که به ترتیب برابر با ۱۳۵ و ($90-180$) روز بود. در خصوص نوع درمان دریافت‌ شده، روش آنژیوپلاستی (۱۰۵ نفر، ۷۰ درصد) شایع‌ترین رویکرد درمانی بوده و پس از آن درمان دارویی (۲۹ نفر، $19.3$ درصد) و جراحی بای‌پس عروق کرونر (۱۶ نفر، $10.7$ درصد) قرار داشتند. بررسی بیماری‌های همراه نشان دهنده شیوع بالای پرفشاری خون (۱۰۵ نفر، ۷۰ درصد) و چربی خون بالا (۱۰۴ نفر، $69.3$ درصد) در میان شرکت‌کنندگان بود؛ سایر بیماری‌های همراه شامل دیابت (۴۴ نفر، $29.3$ درصد)، نارسایی قلبی (۲۶ نفر، $17.3$ درصد)، بیماری‌های ریوی (۱۲ نفر، 8 درصد) و بیماری مزمن کلیوی (۸ نفر، $5.3$ درصد) بودند. ارزیابی سابقه خانوادگی نشان داد که ۱۱۹ نفر ($79.3$ درصد) از بیماران دارای سابقه خانوادگی ابتلا به بیماری‌های قلبی-عروقی بودند.

\section{شاخص‌های پایبندی و رفتارهای سلامت}


تحلیل شاخص‌های پایبندی و رفتارهای سلامت به‌ عنوان اهداف کلیدی پیشگیری ثانویه نشان داد که پایبندی دارویی در ۶۶ درصد از بیماران (۹۹ نفر) در سطح بالا و در ۳۴ درصد (۵۱ نفر) در سطح پایین قرار داشت. هیچ‌یک از بیماران (۱۵۰ نفر، 100 درصد) سابقه شرکت در توانبخشی قلب نداشتند. از منظر سطح سواد سلامت ۸۴ نفر (۵۶ درصد) از افراد سواد سلامت مرزی یا ناکافی داشتند و ۴۴ درصد (۶۶ نفر) از سواد سلامت کافی برخوردار بودند. همچنین ۱۰۴ نفر ($69.3$ درصد) از بیماران درک تهدید پایینی از بیماری خود داشتند و سایرین در گروه‌های درک تهدید متوسط (۴۴ نفر، $29.3$ درصد) و درک تهدید بالا (۲ نفر، $1.3$ درصد) قرار گرفتند.

در بررسی وضعیت مصرف سیگار در 30 روز اخیر، ۱۲۳ نفر (۸۲ درصد) غیرسیگاری و ۲۷ نفر (۱8درصد) سیگاری بودند.

میانه و دامنه میان‌چارکی برای دقایق فعالیت بدنی معادل ۲۴۰ ($120-450$) دقیقه در هفته بود. از نظر پایبندی به رژیم غذایی ، ۶۰ درصد افراد (۹۰ نفر) پایبند به رژیم غذایی بودند. در ارزیابی سلامت روان‌شناختی، وضعیت ۴۶ درصد (۶۹ نفر) از شرکت کنندگان، طبیعی ارزیابی شد. در نهایت، بررسی تجمیعی الگوهای خودمراقبتی مشخص کرد که $17.3$ درصد (۲۶ نفر) از بیماران متعهد به تمام اقدامات پیشگیری ثانویه غیردارویی بودند و سایرین پایبندی ناقصی به این اقدامات داشتند.

{

    \fontsize{12}{14.4}\selectfont % <--- CHANGED: Explicitly sets font size to 12pt

    \renewcommand{\arraystretch}{1.15}

    \setlength{\tabcolsep}{4pt}



    % <--- CHANGED: Replaced table, threeparttable, and tabularx with xltabular

    \begin{xltabular}{\textwidth}{@{} >{\raggedright\arraybackslash}X >{\centering\arraybackslash}X @{}}

        

        % --- 1. First Page Header ---

        \caption{ویژگی‌های دموگرافیک، بالینی و پایبندی به رفتارهای سلامت شرکت‌کنندگان در مطالعه \lr{(n = 150)}}

        \label{tab:descriptive-stats} \\

        \toprule

        \textbf{طبقه‌بندی} & \diagbox[width=6.8cm]{\textbf{فراوانی (درصد)}}{\textbf{میانگین $\pm$ انحراف معیار}} \\

        \midrule

        \endfirsthead



        % --- 2. Subsequent Pages Header ---

        \multicolumn{2}{@{}c@{}}{\tablename~\thetable{} -- ادامه از صفحه قبل} \\

        \toprule

        \textbf{طبقه‌بندی} & \diagbox[width=6.8cm]{\textbf{فراوانی (درصد)}}{\textbf{میانگین $\pm$ انحراف معیار}} \\

        \midrule

        \endhead



        % --- 3. Bottom Footer for Pages that Break ---

        \midrule

        \multicolumn{2}{@{}r@{}}{\textit{ادامه در صفحه بعد...}} \\

        \endfoot



        % --- 4. Final Footer (Includes Tablenotes) ---

        \bottomrule

        \multicolumn{2}{@{}p{\textwidth}@{}}{

            \vspace{6pt}

            \raggedright

            $^{*}$ داده‌ها به‌صورت میانگین $\pm$ انحراف معیار (SD) گزارش شده‌اند.\newline

            $^{**}$ داده‌ها به‌صورت میانه (محدوده بین‌چارکی؛ Q1--Q3) گزارش شده‌اند.\newline

            $^{***}$ سطرهای مربوط به بیماری‌های همراه به‌دلیل وجود بیماری‌های متعدد در یک بیمار، جمع درصد فراوانی آن‌ها بیش از $100$ است.\newline

            \lr{CABG} = \lr{Coronary Artery Bypass Grafting}.

        }

        \endlastfoot



        % ---------------- Demographic characteristics ---------------- %

        \multicolumn{2}{c}{\textbf{ویژگی‌های دموگرافیک}} \\

        \midrule

        سن (سال)$^{*}$ & $62.1 \pm 10.7$ \\

        \cmidrule{1-2}

        مرد        & $105 \; (70.0)$ \\

        زن         & $45 \; (30.0)$ \\

        \cmidrule{1-2}

        متاهل      & $120 \; (80.0)$ \\

        مجرد، مطلقه یا بیوه & $30 \; (20.0)$ \\

        \cmidrule{1-2}

        دیپلم و پایین‌تر      & $119 \; (79.3)$ \\

        تحصیلات دانشگاهی      & $31 \; (20.7)$ \\

        \cmidrule{1-2}

        بیکار، بازنشسته یا خانه‌دار & $84 \; (56.0)$ \\

        شاغل                        & $66 \; (44.0)$ \\

        \cmidrule{1-2}

        شهری   & $89 \; (59.3)$ \\

        روستایی & $61 \; (40.7)$ \\

        \cmidrule{1-2}

        ضعیف        & $65 \; (43.3)$ \\

        متوسط/خوب   & $85 \; (56.7)$ \\

        \cmidrule{1-2}

        بله & $150 \; (100.0)$ \\

        خیر & $0 \; (0.0)$ \\

        \midrule



        % ---------------- Clinical characteristics ---------------- %

        \multicolumn{2}{c}{\textbf{ویژگی‌های بالینی}} \\

        \midrule

        زمان سپری‌شده از سکته قلبی (روز)$^{**}$ & $135 \; (90 - 180)$ \\

        \cmidrule{1-2}

        آنژیوپلاستی              & $105 \; (70.0)$ \\

        درمان دارویی             & $29 \; (19.3)$ \\

        جراحی بای‌پس عروق کرونر (CABG) & $16 \; (10.7)$ \\

        \cmidrule{1-2}

        پرفشاری خون$^{***}$      & $105 \; (70.0)$ \\

        چربی خون بالا$^{***}$   & $104 \; (69.3)$ \\

        دیابت$^{***}$           & $44 \; (29.3)$ \\

        نارسایی قلبی$^{***}$    & $26 \; (17.3)$ \\

        بیماری‌های ریوی$^{***}$  & $12 \; (8.0)$ \\

        بیماری مزمن کلیوی$^{***}$ & $8 \; (5.3)$ \\

        \cmidrule{1-2}

        دارد & $119 \; (79.3)$ \\

        ندارد & $31 \; (20.7)$ \\

        \midrule



        % ---------------- Adherence & health behaviors ---------------- %

        \multicolumn{2}{c}{\textbf{شاخص‌های پایبندی و رفتارهای سلامت}} \\

        \midrule

        پایبندی دارویی --- سطح بالا & $99 \; (66.0)$ \\

        پایبندی دارویی --- سطح پایین & $51 \; (34.0)$ \\

        \cmidrule{1-2}

        سابقه شرکت در توانبخشی قلب --- ندارد & $150 \; (100.0)$ \\

        \cmidrule{1-2}

        سطح سواد سلامت --- مرزی یا ناکافی & $84 \; (56.0)$ \\

        سطح سواد سلامت --- کافی & $66 \; (44.0)$ \\

        \cmidrule{1-2}

        درک تهدید از بیماری --- پایین  & $104 \; (69.3)$ \\

        درک تهدید از بیماری --- متوسط   & $44 \; (29.3)$ \\

        درک تهدید از بیماری --- بالا    & $2 \; (1.3)$ \\

        \cmidrule{1-2}

        وضعیت مصرف سیگار (۳۰ روز اخیر) --- غیرسیگاری & $123 \; (82.0)$ \\

        وضعیت مصرف سیگار (۳۰ روز اخیر) --- سیگاری   & $27 \; (18.0)$ \\

        \cmidrule{1-2}

        فعالیت بدنی (دقیقه در هفته)$^{**}$ & $240 \; (120 - 450)$ \\

        \cmidrule{1-2}

        پایبندی به رژیم غذایی --- بله & $90 \; (60.0)$ \\

        پایبندی به رژیم غذایی --- خیر & $60 \; (40.0)$ \\

        \cmidrule{1-2}

        وضعیت سلامت روان‌شناختی --- طبیعی             & $69 \; (46.0)$ \\

        وضعیت سلامت روان‌شناختی --- غیرطبیعی/نیازمند پیگیری & $81 \; (54.0)$ \\

        \cmidrule{1-2}

        پایبندی به اقدامات پیشگیری ثانویه غیر‌دارویی --- کامل  & $26 \; (17.3)$ \\

        پایبندی به اقدامات پیشگیری ثانویه غیر‌دارویی --- ناقص & $124 \; (82.7)$ \\

        \midrule



        % --- Total --- %

        \multicolumn{1}{@{}r}{\textbf{کل}} & $\mathbf{150 \; (100.0)}$ \\



    \end{xltabular}

}

\section{ارتباط بین متغیر های دموگرافیک و پایبندی به اقدامات پیشگیری ثانویه}

\subsection*{\textbf{قطع مصرف سیگار}}


بر اساس نتایج آزمون تی مستقل سن بیماران ارتباط معکوس و معنی‌داری با وضعیت مصرف سیگار دارد $(p=0.004)$، به طوری که به ازای هر 10 سال افزایش سن (بازه سنی$ 92-31$ سال)، احتمال مصرف سیگار $44.6$ درصد کاهش می‌یابد $(OR=0.94, 95\% \ CI:0.90-0.98)$ .

در بررسی با آزمون مجذور کای، جنسیت بیمار رابطه آماری معنی‌داری با وضعیت مصرف سیگار نشان داد $(p<0.001)$، اما به دلیل عدم وجود فرد سیگاری در میان مبتلایان زن، تعیین نسبت شانس امکان پذیر نبود. همچنین وضعیت اشتغال نیز ارتباط مستقیم و معنی‌داری با این پیامد داشت $(p<0.001)$، به طوری که افراد شاغل شانس بیشتری برای مصرف سیگار در مقایسه با افراد خانه‌دار، بازنشسته یا بیکار داشتند $(OR=10.7, \ 95\% \ CI:3.47-32.94)$. در مقابل، تحلیل‌های آماری نشان داد که وضعیت تاهل، تحصیلات، محل سکونت و وضعیت اقتصادی ارتباط آماری معناداری را نشان ندادند. $(p>0.05)$

\subsection*{\textbf{پایبندی به فعالیت بدنی}}


نتایج آزمون تی مستقل نشان داد که سن بیماران رابطه معکوس و معنی‌داری با پایبندی به فعالیت بدنی دارد $(p<0.001)$، به طوری که با افزایش سن، احتمال پایبندی کاهش می‌یافت $(OR=0.90, \ 95\% \ CI:0.86-0.95)$.

یافته‌های حاصل از آزمون مجذور کای نیز نشان داد که جنسیت با میزان پایبندی به فعالیت بدنی ارتباط معنی‌داری دارد $(p<0.001)$، به طوری که مردان شانس بالاتری برای پایبندی نسبت به زنان داشتند$(OR=0.026, \ 95\% \ CI:0.01-0.07)$. علاوه بر این، وضعیت تاهل نیز ارتباط معنی‌داری با این پیامد نشان داد $(p<0.001)$، به گونه‌ای که شانس پایبندی در افراد متاهل بیشتر از افراد مجرد بود $(OR=4.92, \ 95\% \ CI:2.11-11.46)$. اشتغال نیز رابطه معناداری با فعالیت بدنی داشت $(p<0.001)$. به دلیل عدم وجود فرد غیرپایبند به فعالیت بدنی در میان شاغلین، تعیین نسبت شانس امکان پذیر نبود. همچنین بین سطح تحصیلات و میزان پایبندی به فعالیت بدنی رابطه آماری معنی‌داری یافت شد (\(p=0.049\)). محل سکونت و وضعیت اقتصادی فاقد هرگونه ارتباط آماری معنی‌دار با پایبندی به فعالیت بدنی بودند $(p>0.05)$.

\subsection*{\textbf{پایبندی به رژیم غذایی}}


پیروی از رژیم غذایی مدیترانه‌ای به‌طور معناداری تحت تأثیر وضعیت اقتصادی بود $(p<0.001, \ OR=6.8, \ 95\% \ CI:3.28-14.06)$ و سن، وضعیت تاهل، جنسیت و اشتغال نقش تعیین‌کننده‌ای در آن نداشتند. بیماران دارای تحصیلات دانشگاهی یا عالی، در مقایسه با افراد دارای تحصیلات پایه، شانس به‌مراتب بالاتری را در پیروی از این رژیم غذایی نشان دادند $(p<0.001, \ OR=8.58, \ 95\% \ CI:2.47-29.76)$. همچنین، احتمال پیروی از این الگو در میان ساکنان مناطق شهری بیشتر از روستایی بود $(p<0.001, \ OR=3.9, \ 95\% \ CI:1.95-7.79)$.

\subsection*{\textbf{سلامت روان شناختی}}


یافته‌های آزمون تی مستقل نشان داد که سن رابطه مستقیم و معنی‌داری با میزان دیسترس روان‌شناختی دارد $(p=0.011)$، به طوری که به ازای هر 10 سال سن (بازه سنی 92-31 سال)، احتمال بروز دیسترس روان‌شناختی $50.7$ درصد افزایش می‌یابد $(OR=1.04, \ 95\% \ CI:1.01-1.08)$. تحلیل‌ با آزمون مجذور کای نشان داد که جنسیت بیمار ارتباط معنی‌داری با این پیامد داشت $(p<0.001)$، به طوری که شانس ابتلا به دیسترس روان‌شناختی در زنان 4.48 برابر مردان بود $(OR=4.48, \ 95\% \ CI:2.01-10.00)$. وضعیت تاهل نیز با سلامت روان ارتباط معنی‌داری داشت $(p=0.005)$ و افراد مجرد، مطلقه یا بیوه شانس دیسترس روان‌شناختی بالاتری را در مقایسه با افراد متاهل نشان دادند $(OR=3.51, \ 95\% \ CI:1.40-8.80)$. علاوه بر این، داشتن تحصیلات دانشگاهی $(p=0.003, \ OR=3.74, \ 95\% \ CI:1.58-8.83)$، در کنار وضعیت اشتغال به کار $(p=0.029, \ OR=2.07, \ 95\% \ CI:1.07-3.99)$، به عنوان عوامل محافظتی معنی‌دار در برابر دیسترس روان‌شناختی شناسایی شدند. در مقابل، متغیرهای محل سکونت و وضعیت اقتصادی ارتباط آماری معنی‌داری با وضعیت سلامت روان‌شناختی بیماران نداشتند$(p>0.05)$.

\subsection*{\textbf{پایبندی کلی}}


نتایج آزمون تی مستقل نشان داد که سن بیماران ارتباط آماری معنی‌داری با دستیابی به پایبندی کامل به اقدامات پیشگیری ثانویه ندارد $(p=0.303, \ OR=0.97, \ 95\% \ CI:0.94-1.01)$ بر اساس یافته‌های آزمون مجذور کای، جنسیت بیمار با پایبندی کلی ارتباط معنی‌داری داشت $(p=0.024)$، به طوری که مردان شانس بالاتری برای دستیابی به پایبندی کامل غیردارویی نسبت به زنان بروز دادند $(OR=3.92, \ 95\% \ CI:1.11-13.89)$. سطح تحصیلات نیز قوی‌ترین پیش‌بینی کننده پایبندی کلی بود $(p<0.001, \ OR=7.34, \ 95\% \ CI:2.91-18.52)$. همچنین، وضعیت اقتصادی روندی نزدیک به معنی‌داری با پایبندی کلی نشان داد؛ به طوری که شانس پایبندی کامل در بیماران با وضعیت اقتصادی متوسط/خوب در مقایسه با افراد دارای وضعیت ضعیف، بیش از دو برابر بود، اما این ارتباط از نظر آماری معنی‌دار تلقی نشد $(p=0.063, \ OR=2.38, \ 95\% \ CI:0.93-6.08)$. وضعیت تاهل رابطه معنی‌داری با پایبندی کلی نشان داد $(p=0.024)$؛ با این حال، در تحلیل رگرسیون، ارتباط معنی‌داری را نشان نداد ولی شانس پایبندی کامل در افراد متاهل بالاتر گزارش شد $(p=0.051, \ OR=7.63, \ 95\% \ CI:0.99-58.78)$. در بررسی سایر متغیرها، وضعیت اشتغال و محل سکونت فاقد هرگونه ارتباط آماری معنی‌دار با پیامد پایبندی کلی غیردارویی در بیماران پس از انفارکتوس میوکارد بودند $(p>0.05)$.

\section{ارتباط بین متغیر های بالینی و پایبندی به اقدامات پیشگیری ثانویه}

\subsection*{\textbf{قطع مصرف سیگار}}


در بررسی ها، سواد سلامت رابطه معنی‌داری با مصرف سیگار نشان داد $(p=0.028)$، به گونه‌ای که شانس مصرف سیگار در افراد با سواد سلامت کافی $2.57$ برابر بیشتر از گروه مرزی یا نامناسب بود (فاصله اطمینان: $1.09$ تا $6.07$). با این حال، این رابطه به شدت تحت تاثیر سن به عنوان یک متغیر مخدوش‌کننده قرار داشت؛ از آنجا که بیماران با سواد سلامت کافی، میانگین سنی به مراتب پایین‌تری داشتند ($54.14$ در برابر $68.33$سال)، پس از تعدیل اثر سن در مدل رگرسیون لجستیک، این ارتباط متناقض کاملاً ناپدید و غیرمعنی‌دار شد $(p=0.740)$، نسبت شانس تعدیل‌شده: $1.22$). در مقابل، تعداد روزهای سپری‌شده از سکته، نوع درمان، تعداد بیماری‌های زمینه‌ای، سابقه خانوادگی، درک کل از بیماری و میزان تبعیت دارویی ارتباط آماری معناداری با این پیامد نداشتند $(p>0.05)$.

\subsection*{\textbf{پایبندی به فعالیت بدنی}}


نتایج آزمون من-ویتنی یو نشان دهنده عدم وجود ارتباط معنی‌دار بین تعداد روزهای سپری‌شده از سکته قلبی و پایبندی به فعالیت بدنی بود $(p=0.799)$ و تحلیل رگرسیون لجستیک نیز این عدم ارتباط را تایید کرد 
$(OR=1.00, \ 95\% \ CI:0.99-1.01)$. در بررسی متغیر نوع درمان، آزمون مجذور کای انجام شد که نشان‌دهنده یک ارتباط آماری معنی‌دار بود $(p<0.001)$. به طوری که با در نظر گرفتن روش آنژیوپلاستی به عنوان گروه مرجع در رگرسیون لجستیک، بیماران تحت درمان دارویی با احتمال کمتری به فعالیت بدنی پایبند بودند $(OR=0.15, \ 95\% \ CI:0.06-0.36)$، در حالی که تفاوت معنی‌دار آماری میان بیماران تحت عمل جراحی بای‌پاس عروق کرونر و گروه مرجع مشاهده نشد $(OR=0.84, \ 95\% \ CI:0.25-2.86)$. همچنین میان تعداد بیماری‌های زمینه‌ای و فعالیت فیزیکی ارتباط آماری معنی‌داری دیده شد $(p=0.019)$. به گونه‌ای که به ازای هر یک بیماری زمینه‌ای، شانس پایبندی به فعالیت بدنی کاهش پیدا کرد $(OR=0.63, \ 95\% \ CI:0.46-0.87)$.

متغیر سابقه خانوادگی بیماری‌های قلبی-عروقی ارتباط معنی‌داری با میزان پایبندی به فعالیت بدنی نشان نداد $(p=0.660, \ OR=0.74, \ 95\% \ CI:0.30-1.81)$. در مقابل، وضعیت سواد سلامت رابطه معنی‌داری با میزان پایبندی به فعالیت بدنی داشت $(p<0.001)$. به طوری که شانس پایبندی به فعالیت بدنی در بیماران دارای سواد سلامت کافی، حدود 5 برابر بیشتر از بیماران با سواد سلامت مرزی یا ناکافی بود $(OR=4.99, \ 95\% \ CI:2.19-11.36)$ افزون بر این، توزیع نمره کل پرسشنامه درک از بیماری بین دو گروه پایبند و غیرپایبند به فعالیت بدنی تفاوت معنی‌داری دارد $(p=0.025)$. با این حال، در تحلیل رگرسیون لوجستیک تک‌متغیره، نمره کل \lr{IPQ} ارتباط معنی‌داری با پایبندی به فعالیت بدنی نداشت $(p =0.058, \ OR=0.916, \ 95\% \ CI:0.836-1.003)$. با این حال، ارتباط میان پایبندی دارویی و پایبندی به فعالیت بدنی معنادار بود $(p=0.003)$. به گونه‌ای که بیماران با پایبندی دارویی ضعیف، با احتمال کمتری از توصیه‌های فعالیت بدنی پیروی می کردند $(OR=0.32, \ 95\% \ CI:0.16-0.66)$.


\subsection*{\textbf{پایبندی به رژیم غذایی}}


پایبندی به رژیم غذایی مدیترانه‌ای ارتباط آماری معناداری را با دو متغیر سواد سلامت و پایبندی دارویی نشان داد؛ به‌طوری که شانس پایبندی به این الگوی غذایی در بیماران دارای سواد سلامت محدود یا ناکافی، کاهش ۶۷ درصدی را نسبت به افراد با سواد سلامت کافی نشان داد $(p =0.003, \ OR=0.33, \ 95\% \ CI:0.16-0.66)$. همچنین، بیماران با پایبندی دارویی بالا، به‌طور میانگین $2.5$ برابر شانس بیشتری برای تبعیت از رژیم غذایی مدیترانه‌ای در مقایسه با گروه با پایبندی دارویی کم داشتند $(p =0.012, \ OR=2.55, \ 95\% \ CI:1.27-5.10)$. در مقابل، متغیرهای زمان سپری‌شده از سکته قلبی، نمره کل درک از بیماری، تعداد بیماری‌های زمینه‌ای، نوع مداخله درمانی و سابقه خانوادگی بیماری‌های قلبی-عروقی ارتباط آماری معناداری با میزان پایبندی به رژیم غذایی نداشتند $(p>0.05)$.


\subsection*{\textbf{سلامت روان شناختی}}


نتایج بررسی ها نشان داد که درک بیمار از بیماری خود رابطه مستقیم و معنی‌داری با میزان دیسترس روان‌شناختی دارد $(p<0.001)$، به طوری که میانگین نمرات در گروه با دیسترس روان‌شناختی بالاتر از گروه نرمال بود $(p < 0.001, \ OR=1.31, \ 95\% \ CI:1.16 - 1.48)$. علاوه بر این، بررسی‌های دومتغیره نشان داد که سطح سواد سلامت نیز ارتباط آماری معنی‌داری با دیسترس روان‌شناختی دارد $(p=0.007)$ و شانس بروز دیسترس روان‌شناختی در بیماران با سواد سلامت مرزی/ناکافی $2.6$ برابر بیشتر از بیماران با سواد سلامت کافی بود $(OR=2.60, \ 95\% \ CI:1.34-5.05)$. در مقابل، متغیرهای تعداد روزهای سپری‌شده از انفارکتوس میوکارد، تعداد بیماری‌های زمینه‌ای، نوع درمان دریافتی، سابقه خانوادگی و میزان پیروی از درمان‌های دارویی ارتباط آماری معنی‌داری با وضعیت دیسترس روان‌شناختی بیماران نداشتند $(p>0.05)$.

\subsection*{\textbf{پایبندی کلی}}


سطح سواد سلامت بیمار ارتباط معنی‌داری با پایبندی کامل به درمان نشان داد (\(p=0.016\))، به طوری که بیماران دارای سواد سلامت کافی شانس بالاتری را برای دستیابی به پایبندی کامل در مقایسه با افراد دارای سواد سلامت مرزی یا ناکافی نشان دادند $(OR=2.891, \ 95\% \ CI:1.194-7.002)$. همچنین، وضعیت پایبندی دارویی ارتباط آماری معنی‌داری با پایبندی کلی غیردارویی داشت $(p=0.008)$ ؛ به گونه‌ای که شانس پایبندی کلی در بیماران با پایبندی دارویی بالا، بیشتر از افراد با پایبندی دارویی پایین گزارش شد $(OR=4.842, \ 95\% \ CI:1.379-17.006)$. علاوه بر این، ارزیابی نمره کل درک از بیماری نشان‌دهنده تفاوت معنی‌دار میان دو گروه پایبند و غیرپایبند بود $(p=0.007, \ OR=0.873, 95\% \ CI:0.773-0.986)$. در بررسی سایر متغیرهای مستقل، تعداد روزهای سپری‌شده از انفارکتوس میوکارد، تعداد بیماری‌های زمینه‌ای، نوع درمان دریافتی و سابقه خانوادگی فاقد ارتباط آماری معنی‌دار با پایبندی کلی در بیماران بودند $(p>0.05)$.


\begin{figure}[H]
  \centering
  \begin{tikzpicture}
    \begin{axis}[
      width=0.98\textwidth,
      height=0.68\textwidth,
      xmin=-0.70,
      xmax=7.70,
      ymin=0,
      ymax=108,
      ybar,
      bar width=18pt,
      axis background/.style={fill=Canvas},
      axis x line*=bottom,
      axis y line*=left,
      axis line style={draw=Ink, line width=0.7pt},
      tick style={draw=none},
      ymajorgrids=true,
      grid style={draw=Border, line width=0.45pt},
      ytick={0,20,40,60,80,100},
      yticklabel style={font=\small, text=Ink},
      ylabel={درصد},
      ylabel style={font=\normalsize\bfseries, text=Ink, yshift=-2pt},
      xtick={0,1,3,4,6,7},
      xticklabels={
        {سطح ناکافی},
        {سطح کافی},
        {غیر شاغل},
        {شاغل},
        {زن},
        {مرد}
        },
      xticklabel style={font=\small, text=Ink, align=center, yshift=-2pt},
      enlarge x limits=false,
      clip=false,
      title={\bfseries عدم مصرف سیگار (\lr{n=123})},
      title style={font=\Large, text=Ink, yshift=37pt},
      legend columns=2,
      legend style={
        at={(0.5,1.02)},
        anchor=south,
        draw=Border,
        fill=SurfaceOne,
        rounded corners=2pt,
        inner xsep=8pt,
        inner ysep=4pt,
        font=\small,
        text=Ink,
        /tikz/every even column/.append style={column sep=14pt}
      },
      legend cell align={right}
    ]
      % In the right-to-left layout, each first category is the purple bar
      % on the right side of its pair.
      \addplot+[
        draw=PrimaryPurple,
        fill=PrimaryPurple,
        nodes near coords,
        point meta=explicit symbolic,
        every node near coord/.append style={
          font=\footnotesize\bfseries,
          text=Ink,
          anchor=south,
          yshift=1.5pt
        }
      ] coordinates {
        (1,74.6) [74.6٪]
        (4,66.2) [66.2٪]
        (7,74.3) [74.3٪]
      };
      \addlegendentry{ردهٔ نخست هر گروه}
      \addplot+[
        draw=SecondaryTeal,
        fill=SecondaryTeal,
        nodes near coords,
        point meta=explicit symbolic,
        every node near coord/.append style={
          font=\footnotesize\bfseries,
          text=Ink,
          anchor=south,
          yshift=1.5pt
        }
      ] coordinates {
        (0,88) [88٪]
        (3,95.1) [95.1٪]
        (6,100) [100٪]
      };
      \addlegendentry{ردهٔ دوم هر گروه}
      % Group names and exact p-values, centered below the category labels.
      \node[align=center, font=\small\bfseries, text=Ink]
        at (axis cs:0.5,-17) {سواد سلامت\\[-1pt]
        \textcolor{InkMuted}{\normalfont\lr{$p=0.035$}}};
      \node[align=center, font=\small\bfseries, text=Ink]
        at (axis cs:3.5,-17) {اشتغال\\[-1pt]
        \textcolor{InkMuted}{\normalfont\lr{$p<0.001$}}};
      \node[align=center, font=\small\bfseries, text=Ink]
        at (axis cs:6.5,-17) {جنسیت\\[-1pt]
        \textcolor{InkMuted}{\normalfont\lr{$p<0.001$}}};
    \end{axis}
  \end{tikzpicture}
  \caption{درصد عدم مصرف سیگار به تفکیک جنسیت، وضعیت اشتغال و سطح سواد سلامت. مقادیر روی ستون‌ها درصد را نشان می‌دهند و مقدارهای \lr{p} زیر هر گروه گزارش شده‌اند.}
  \label{fig:no-smoking-grouped-bar}
\end{figure}

\begin{figure}[H]
    \centering
    % TikZ is placed in an LTR environment to keep its coordinate system stable.
    % Individual Persian text elements are restored with \faText.
    \begin{tikzpicture}
        \begin{axis}[
            width=0.99\textwidth,
            height=10cm,
            % Axis ranges
            ymin=0,
            ymax=112,
            xmin=0.15,
            xmax=21.85,
            % Background and axis lines
            axis background/.style={
                fill=Canvas
            },
            axis x line*=bottom,
            axis y line*=left,
            axis line style={
                draw=Ink,
                line width=0.75pt
            },
            tick style={
                draw=InkMuted,
                line width=0.55pt
            },
            % Y-axis
            ytick={0,20,40,60,80,100},
            yticklabels={
                \faText{۰},
                \faText{۲۰},
                \faText{۴۰},
                \faText{۶۰},
                \faText{۸۰},
                \faText{۱۰۰}
            },
            ylabel={\faText{درصد فراوانی}},
            ylabel style={
                font=\bfseries\small,
                text=Ink,
                yshift=-2pt
            },
            yticklabel style={
                font=\small,
                text=InkMuted
            },
            % Grid
            xmajorgrids=false,
            ymajorgrids=true,
            grid style={
                draw=Border,
                line width=0.55pt
            },
            % One tick beneath every individual bar
            xtick={
                1,2,
                4,5,
                7,8,
                10,11,
                13,14,15,
                17,18,
                20,21
            },
            xticklabels={
                \faText{مرد},
                \faText{زن},
                \faText{مجرد},
                \faText{متاهل},
                \faText{شاغل},
                \faText{غیر شاغل},
                \faText{پایه},
                \faText{دانشگاهی},
                \faText{آنژیوپلاستی},
                \faText{جراحی},
                \faText{دارویی},
                \faText{سطح کافی},
                \faText{سطح ناکافی},
                \faText{بالا},
                \faText{پایین}
            },
            xticklabel style={
                font=\footnotesize,
                text=Ink,
                rotate=52,
                anchor=east,
                align=right,
                yshift=-2pt
            },
            tick align=outside,
            % Bars and exact percentage labels
            bar width=13pt,
            clip=false,
            point meta=explicit symbolic,
            nodes near coords={\pgfplotspointmeta},
            every node near coord/.append style={
                font=\bfseries\scriptsize,
                text=Ink,
                anchor=south,
                yshift=2pt,
                inner sep=0.5pt
            },
            % Chart header
]
        %--------------------------------------------------------------
        % First category of every independent variable
        %--------------------------------------------------------------
        \addplot+[
            ybar,
            mark=none,
            draw=PrimaryPurple,
            fill=PrimaryPurple
        ] coordinates {
            (1,90.5) [{۹۰٫۵٪}]
            (4,40) [{۴۰٪}]
            (7,100) [{۱۰۰٪}]
            (10,65.5) [{۶۵٫۵٪}]
            (13,78.1) [{۷۸٫۱٪}]
            (17,86.6) [{۸۶٫۶٪}]
            (20,77.8) [{۷۷٫۸٪}]
        };
        %--------------------------------------------------------------
        % Second category of every independent variable
        %--------------------------------------------------------------
        \addplot+[
            ybar,
            mark=none,
            draw=SecondaryTeal,
            fill=SecondaryTeal
        ] coordinates {
            (2,20) [{۲۰٪}]
            (5,76.7) [{۷۶٫۷٪}]
            (8,45.2) [{۴۵٫۲٪}]
            (11,83.9) [{۸۳٫۹٪}]
            (14,75) [{۷۵٪}]
            (18,55.4) [{۵۵٫۴٪}]
            (21,52.9) [{۵۲٫۹٪}]
        };
        %--------------------------------------------------------------
        % Third category, used only for the three-level treatment group
        %--------------------------------------------------------------
        \addplot+[
            ybar,
            mark=none,
            draw=AccentMagenta,
            fill=AccentMagenta
        ] coordinates {
            (15,34.5) [{۳۴٫۵٪}]
        };
        %--------------------------------------------------------------
        % Independent-variable names and p-values
        %--------------------------------------------------------------
        \node[
            font=\bfseries\scriptsize,
            text=Ink,
            align=center,
            text width=1.9cm
        ]
        at (axis cs:1.5,-28) {
            \faText{جنسیت}\\[-1pt]
            {$p<0.001$}
        };
        \node[
            font=\bfseries\scriptsize,
            text=Ink,
            align=center,
            text width=1.9cm
        ]
        at (axis cs:4.5,-28) {
            \faText{تاهل}\\[-1pt]
            {$p<0.001$}
        };
        \node[
            font=\bfseries\scriptsize,
            text=Ink,
            align=center,
            text width=1.9cm
        ]
        at (axis cs:7.5,-28) {
            \faText{اشتغال}\\[-1pt]
            {$p<0.001$}
        };
        \node[
            font=\bfseries\scriptsize,
            text=Ink,
            align=center,
            text width=1.9cm
        ]
        at (axis cs:10.5,-28) {
            \faText{تحصیلات}\\[-1pt]
            {$p=0.049$}
        };
        \node[
            font=\bfseries\scriptsize,
            text=Ink,
            align=center,
            text width=1.9cm
        ]
        at (axis cs:14,-28) {
            \faText{نوع درمان}\\[-1pt]
            {$p<0.001$}
        };
        \node[
            font=\bfseries\scriptsize,
            text=Ink,
            align=center,
            text width=1.9cm
        ]
        at (axis cs:17.5,-28) {
            \faText{سواد سلامت}\\[-1pt]
            {$p<0.001$}
        };
        \node[
            font=\bfseries\scriptsize,
            text=Ink,
            align=center,
            text width=1.9cm
        ]
        at (axis cs:20.5,-28) {
            \faText{پایبندی دارویی}\\[-1pt]
            {$p=0.002$}
        };
        \end{axis}
    \end{tikzpicture}
    \caption{
        درصد فراوانی پایبندی به فعالیت بدنی کافی}
    \label{fig:sufficient-physical-activity}
\end{figure}

\begin{figure}[H]
    \centering
    % TikZ is placed in an LTR environment to keep its coordinate system stable.
    % Individual Persian text elements are restored with \faText.
    \begin{tikzpicture}
        \begin{axis}[
            width=0.85\textwidth,
            height=9cm,
            % Axis ranges
            ymin=0,
            ymax=112,
            xmin=0.15,
            xmax=14.85,
            % Background and axis lines
            axis background/.style={
                fill=Canvas
            },
            axis x line*=bottom,
            axis y line*=left,
            axis line style={
                draw=Ink,
                line width=0.75pt
            },
            tick style={
                draw=InkMuted,
                line width=0.55pt
            },
            % Y-axis
            ytick={0,20,40,60,80,100},
            yticklabels={
                \faText{۰},
                \faText{۲۰},
                \faText{۴۰},
                \faText{۶۰},
                \faText{۸۰},
                \faText{۱۰۰}
            },
            ylabel={\faText{درصد فراوانی}},
            ylabel style={
                font=\bfseries\small,
                text=Ink,
                yshift=-2pt
            },
            yticklabel style={
                font=\small,
                text=InkMuted
            },
            % Grid
            xmajorgrids=false,
            ymajorgrids=true,
            grid style={
                draw=Border,
                line width=0.55pt
            },
            % One tick beneath every individual bar
            xtick={
                1,2,
                4,5,
                7,8,
                10,11,
                13,14
            },
            xticklabels={
                \faText{پایه},
                \faText{دانشگاهی},
                \faText{روستایی},
                \faText{شهری},
                \faText{ضعیف},
                \faText{متوسط/خوب},
                \faText{سطح کافی},
                \faText{سطح ناکافی},
                \faText{بالا},
                \faText{پایین}
            },
            xticklabel style={
                font=\footnotesize,
                text=Ink,
                rotate=52,
                anchor=east,
                align=right,
                yshift=-2pt
            },
            tick align=outside,
            % Bars and exact percentage labels
            bar width=13pt,
            clip=false,
            point meta=explicit symbolic,
            nodes near coords={\pgfplotspointmeta},
            every node near coord/.append style={
                font=\bfseries\scriptsize,
                text=Ink,
                anchor=south,
                yshift=2pt,
                inner sep=0.5pt
            },
]
        %--------------------------------------------------------------
        % First category of every independent variable
        %--------------------------------------------------------------
        \addplot+[
            ybar,
            mark=none,
            draw=PrimaryPurple,
            fill=PrimaryPurple
        ] coordinates {
            (1,52.1) [{۵۲٫۱٪}]
            (4,41.0) [{۴۱٪}]
            (7,35.4) [{۳۵٫۴٪}]
            (10,74.6) [{۷۴٫۶٪}]
            (13,67.3) [{۶۷٫۳٪}]
        };

        %--------------------------------------------------------------
        % Second category of every independent variable
        %--------------------------------------------------------------
        \addplot+[
            ybar,
            mark=none,
            draw=SecondaryTeal,
            fill=SecondaryTeal
        ] coordinates {
            (2,90.3) [{۹۰٫۳٪}]
            (5,73.0) [{۷۳٪}]
            (8,78.8) [{۷۸٫۸٪}]
            (11,48.2) [{۴۸٫۲٪}]
            (14,44.9) [{۴۴٫۹٪}]
        };

        %--------------------------------------------------------------
        % Independent-variable names and p-values
        %--------------------------------------------------------------
        \node[
            font=\bfseries\scriptsize,
            text=Ink,
            align=center,
            text width=1.9cm
        ]
        at (axis cs:1.5,-29) {
            \faText{تحصیلات}\\[-1pt]
            {$p<0.001$}
        };

        \node[
            font=\bfseries\scriptsize,
            text=Ink,
            align=center,
            text width=1.9cm
        ]
        at (axis cs:4.5,-29) {
            \faText{سکونت}\\[-1pt]
            {$p<0.001$}
        };

        \node[
            font=\bfseries\scriptsize,
            text=Ink,
            align=center,
            text width=1.9cm
        ]
        at (axis cs:7.5,-29) {
            \faText{وضعیت اقتصادی}\\[-1pt]
            {$p<0.001$}
        };

        \node[
            font=\bfseries\scriptsize,
            text=Ink,
            align=center,
            text width=1.9cm
        ]
        at (axis cs:10.5,-29) {
            \faText{سواد سلامت}\\[-1pt]
            {$p=0.001$}
        };

        \node[
            font=\bfseries\scriptsize,
            text=Ink,
            align=center,
            text width=1.9cm
        ]
        at (axis cs:13.5,-29) {
            \faText{پایبندی دارویی}\\[-1pt]
            {$p=0.009$}
        };

        \end{axis}
    \end{tikzpicture}
    \caption{درصد فراوانی پایبندی به رژیم غذایی}
    \label{fig:diet-adherence-grouped-bar}
\end{figure}

\begin{figure}[H]
    \centering
    % TikZ is placed in an LTR coordinate system.
    % Individual Persian text elements are restored with \faText.
    \begin{tikzpicture}
        \begin{axis}[
            width=0.85\textwidth,
            height=8.5cm,
            % Axis ranges
            ymin=0,
            ymax=88,
            xmin=0.15,
            xmax=14.85,
            % Background and axis lines
            axis background/.style={
                fill=Canvas
            },
            axis x line*=bottom,
            axis y line*=left,
            axis line style={
                draw=Ink,
                line width=0.75pt
            },
            tick style={
                draw=InkMuted,
                line width=0.55pt
            },
            % Y-axis
            ytick={0,20,40,60,80},
            yticklabels={
                \faText{۰},
                \faText{۲۰},
                \faText{۴۰},
                \faText{۶۰},
                \faText{۸۰}
            },
            ylabel={\faText{درصد فراوانی}},
            ylabel style={
                font=\bfseries\small,
                text=Ink,
                yshift=-2pt
            },
            yticklabel style={
                font=\small,
                text=InkMuted
            },
            % Grid
            xmajorgrids=false,
            ymajorgrids=true,
            grid style={
                draw=Border,
                line width=0.55pt
            },
            % One tick beneath every individual bar
            xtick={
                1,2,
                4,5,
                7,8,
                10,11,
                13,14
            },
            xticklabels={
                \faText{مرد},
                \faText{زن},
                \faText{مجرد},
                \faText{متاهل},
                \faText{شاغل},
                \faText{غیر شاغل},
                \faText{پایه},
                \faText{دانشگاهی},
                \faText{سطح کافی},
                \faText{سطح ناکافی}
            },
            xticklabel style={
                font=\footnotesize,
                text=Ink,
                rotate=52,
                anchor=east,
                align=right,
                yshift=-2pt
            },
            tick align=outside,
            % Bars and exact percentage labels
            bar width=13pt,
            clip=false,
            point meta=explicit symbolic,
            nodes near coords={\pgfplotspointmeta},
            every node near coord/.append style={
                font=\bfseries\scriptsize,
                text=Ink,
                anchor=south,
                yshift=2pt,
                inner sep=0.5pt
            }
        ]

        %--------------------------------------------------------------
        % First category of every independent variable
        %--------------------------------------------------------------
        \addplot+[
            ybar,
            mark=none,
            draw=PrimaryPurple,
            fill=PrimaryPurple
        ] coordinates {
            (1,56.2) [{۵۶٫۲٪}]
            (4,23.3) [{۲۳٫۳٪}]
            (7,56)   [{۵۶٪}]
            (10,39.5) [{۳۹٫۵٪}]
            (13,59.1) [{۵۹٫۱٪}]
        };

        %--------------------------------------------------------------
        % Second category of every independent variable
        %--------------------------------------------------------------
        \addplot+[
            ybar,
            mark=none,
            draw=SecondaryTeal,
            fill=SecondaryTeal
        ] coordinates {
            (2,22.2) [{۲۲٫۲٪}]
            (5,51.7) [{۵۱٫۷٪}]
            (8,38.1) [{۳۸٫۱٪}]
            (11,71)  [{۷۱٪}]
            (14,35.7) [{۳۵٫۷٪}]
        };

        %--------------------------------------------------------------
        % Independent-variable names and p-values
        %--------------------------------------------------------------
        \node[
            font=\bfseries\scriptsize,
            text=Ink,
            align=center,
            text width=1.9cm
        ]
        at (axis cs:1.5,-27) {
            \faText{جنسیت}\\[-1pt]
            \normalfont {$p<0.001$}
        };

        \node[
            font=\bfseries\scriptsize,
            text=Ink,
            align=center,
            text width=1.9cm
        ]
        at (axis cs:4.5,-27) {
            \faText{تاهل}\\[-1pt]
            \normalfont {$p=0.005$}
        };

        \node[
            font=\bfseries\scriptsize,
            text=Ink,
            align=center,
            text width=1.9cm
        ]
        at (axis cs:7.5,-27) {
            \faText{اشتغال}\\[-1pt]
            \normalfont {$p=0.029$}
        };

        \node[
            font=\bfseries\scriptsize,
            text=Ink,
            align=center,
            text width=1.9cm
        ]
        at (axis cs:10.5,-27) {
            \faText{تحصیلات}\\[-1pt]
            \normalfont {$p=$0.003}
        };

        \node[
            font=\bfseries\scriptsize,
            text=Ink,
            align=center,
            text width=1.9cm
        ]
        at (axis cs:13.5,-27) {
            \faText{سواد سلامت}\\[-1pt]
            \normalfont {$p=0.007$}
        };

        \end{axis}
    \end{tikzpicture}

    \caption{درصد فراوانی سلامت روانشناختی}
    \label{fig:psychological-health-grouped-bar-chart}
\end{figure}

\begin{figure}[H]
    \centering
    % The LTR environment prevents RTL transformations
    % from affecting the PGFPlots coordinate system.
    \begin{latin}
    \begin{tikzpicture}
        \begin{axis}[
            width=0.98\textwidth,
            height=12.2cm,
            ybar,
            bar width=13pt,
            ymin=0,
            ymax=105,
            ytick={0,20,40,60,80,100},
            ylabel={\rl{درصد}},
            ylabel style={
                font=\bfseries,
                text=Ink,
                yshift=-2pt
            },
            title={\rl{\bfseries پایبندی کلی} \lr{(n=26)}},
            title style={
                font=\Large,
                text=Ink,
                yshift=7pt
            },
            axis background/.style={
                fill=Canvas
            },
            axis x line*=bottom,
            axis y line*=left,
            axis line style={
                draw=Ink,
                line width=0.8pt
            },
            tick style={
                draw=Ink,
                line width=0.7pt
            },
            ymajorgrids=true,
            xmajorgrids=false,
            major grid style={
                draw=Border,
                line width=0.65pt
            },
            xtick={
                1,2,
                4,5,
                7,8,
                10,11,
                13,14
            },
            xticklabels={
                \rl{مرد},
                \rl{زن},
                \rl{مجرد},
                \rl{متاهل},
                \rl{پایه},
                \rl{دانشگاهی},
                \rl{سطح کافی},
                \rl{سطح ناکافی},
                \rl{بالا},
                \rl{پایین}
            },
            x tick label style={
                font=\scriptsize,
                text=Ink,
                align=right,
                rotate=35,
                anchor=north east,
                yshift=-1pt
            },
            y tick label style={
                font=\small,
                text=InkMuted
            },
            enlarge x limits=0.065,
            clip=false,
            legend columns=1,
            legend style={
                at={(0.98,0.97)},
                anchor=north east,
                draw=Border,
                fill=SurfaceOne,
                rounded corners=2pt,
                inner xsep=8pt,
                inner ysep=5pt,
                row sep=2pt,
                font=\small,
                text=Ink
            },
            nodes near coords,
            point meta=explicit symbolic,
            every node near coord/.append style={
                font=\small,
                text=Ink,
                anchor=south,
                yshift=2pt
            }
        ]
        % First category in every pair:
        % primary purple, #7B2FBE.
        \addplot+[
            draw=PrimaryPurple,
            fill=PrimaryPurple,
            line width=0.8pt
        ] coordinates {
            (1,21.9)  [\PercentLabel{21.9}]
            (4,3.3)   [\PercentLabel{3.3}]
            (7,10.1)  [\PercentLabel{10.1}]
            (10,26.1) [\PercentLabel{26.1}]
            (13,23.2) [\PercentLabel{23.2}]
        };
        \addlegendentry{\rl{دسته نخست هر جفت}}
        % Second category in every pair:
        % secondary teal, #00C4CC.
        \addplot+[
            draw=SecondaryTeal,
            fill=SecondaryTeal,
            line width=0.8pt
        ] coordinates {
            (2,6.7)   [\PercentLabel{6.7}]
            (5,20.8)  [\PercentLabel{20.8}]
            (8,45.2)  [\PercentLabel{45.2}]
            (11,10.6) [\PercentLabel{10.6}]
            (14,5.9)  [\PercentLabel{5.9}]
        };
        \addlegendentry{\rl{دسته دوم هر جفت}}
        % Variable names and p-values centered
        % below the corresponding pairs of bars.
        \node[
            font=\small,
            text=Ink,
            align=center,
            anchor=north
        ] at (axis cs:1.5,-18) {
            \GenderGroupLabel
        };
        \node[
            font=\small,
            text=Ink,
            align=center,
            anchor=north
        ] at (axis cs:4.5,-18) {
            \MaritalGroupLabel
        };
        \node[
            font=\small,
            text=Ink,
            align=center,
            anchor=north
        ] at (axis cs:7.5,-18) {
            \EducationGroupLabel
        };
        \node[
            font=\small,
            text=Ink,
            align=center,
            anchor=north
        ] at (axis cs:10.5,-18) {
            \LiteracyGroupLabel
        };
        \node[
            font=\small,
            text=Ink,
            align=center,
            anchor=north
        ] at (axis cs:13.5,-18) {
            \MedicationGroupLabel
        };
        \end{axis}
    \end{tikzpicture}
    \end{latin}
    \caption{مقایسه درصد پایبندی کلی برحسب جنسیت، وضعیت تاهل، سطح تحصیلات، سواد سلامت و پایبندی دارویی}
    \label{fig:overall-adherence-grouped-bar}
\end{figure}


\section{تحلیل چندمتغیره عوامل مرتبط با پایبندی به اقدامات پیشگیری ثانویه}

\subsection*{\textbf{قطع مصرف سیگار}}


یافته‌های حاصل از تحلیل چندمتغیره\footnote{\lr{Firth Penalized Logistic Regression}} نشان داد که جنسیت مذکر ارتباط مستقیم و معنی‌داری با افزایش شانس مصرف سیگار دارد. به طوری که شانس مصرف سیگار در مردان بالاتر از زنان برآورد شد $(p=0.017)$. علاوه بر این، وضعیت اشتغال نیز رابطه معنی‌داری را نشان داد و افراد بیکار، بازنشسته یا خانه‌دار در مقایسه با افراد شاغل، با احتمال کمتری سیگار مصرف می کردند \\$(p=0.037)$.

در مقابل، سن و سواد سلامت ارتباط آماری معنی‌داری را با وضعیت مصرف سیگار در بیماران مورد مطالعه نشان ندادند $(p>0.05)$.



\subsection*{\textbf{پایبندی به فعالیت بدنی}}


نتایج بررسی ها نشان داد که جنسیت ارتباط معناداری با پایبندی به فعالیت بدنی دارد، به طوری که مردان شانس بالاتری برای تبعیت از برنامه‌های فعالیت بدنی نسبت به زنان دارند $(p<0.001)$. وضعیت اشتغال نیز به عنوان یک عامل مرتبط شناخته شد و گروه بیکار، بازنشسته یا خانه‌دار در مقایسه با افراد شاغل پایبندی کمتری داشتند $(p=0.001)$. همچنین بیمارانی که تحت درمان دارویی قرار داشتند، در مقایسه با بیماران تحت درمان با آنژیوپلاستی، به شکل معنی‌داری پایبندی کمتری به فعالیت بدنی داشتند $(p=0.031)$ ولی جراحی بای‌پاس عروق کرونر، در مقایسه با آنژیوپلاستی، تفاوت معنی داری در فعالیت بدنی نداشتند $(p>0.05)$.

از سوی دیگر، سن، وضعیت تاهل، سطح تحصیلات، تعداد بیماری‌های همراه، سواد سلامت، نمره پرسشنامه درک از بیماری و پایبندی دارویی ارتباط معنی‌داری با میزان پایبندی به فعالیت بدنی نشان ندادند $(p>0.05)$.




\subsection*{\textbf{پایبندی به رژیم غذایی}}


از میان تمامی متغیرها، تنها وضعیت اقتصادی ارتباط مستقیم و معنی‌داری با پایبندی به رژیم غذایی داشت. به طوری که وضعیت اقتصادی بهتر به عنوان یک عامل پیش‌بینی کننده، شانس پایبندی به رژیم غذایی را در بیماران پس از انفارکتوس میوکارد افزایش می‌دهد $(p=0.002)$.
در مقابل، سطح تحصیلات بیماران، محل سکونت، سواد سلامت و پایبندی دارویی رابطه معنی‌دار آماری را با وضعیت پایبندی به رژیم غذایی نشان ندادند $(p>0.05)$.



\subsection*{\textbf{سلامت روان شناختی}}


متغیرهای وضعیت تاهل، جنسیت و ادراک بیماری ارتباط مستقیم و معنی‌دار آماری با سلامت روان‌شناختی داشتند. بررسی‌ها نشان داد که بیماران متاهل با احتمال بالاتری سلامت روانشناختی داشتند $(p=0.047)$. از سوی دیگر، زنان به طور قابل‌توجهی سلامت روانشناختی نا‌مطلوب تری داشتند $(p=0.005)$. افزون بر این، ذهنیت منفی‌تر نسبت به بیماری با کاهش سطح سلامت روانشناختی همراه بود $(p=0.001)$.

سایر متغیرهای وارد شده به این تحلیل از جمله سن، سطح تحصیلات، وضعیت اشتغال بیماران و سطح سواد سلامت فاقد ارتباط آماری معنی‌دار با سلامت روانشناختی بودند $(p>0.05)$.


\subsection*{\textbf{پایبندی کلی}}


تحلیل آماری نشان داد که تنها سطح تحصیلات بیماران ارتباط مستقیم و معنی‌داری با پایبندی کامل به اقدامات پیشگیری ثانویه پس از انفارکتوس میوکارد دارد. به طوری که بیماران دارای تحصیلات دانشگاهی در مقایسه با افراد دارای سطح تحصیلات دیپلم و پایین‌تر، پیروی بیشتری از توصیه‌های پیشگیرانه نشان دادند $(p=0.003)$.

در مقابل، وضعیت تأهل بیماران، جنسیت، سواد سلامت، درک از بیماری و میزان پایبندی دارویی، رابطه معنی‌دار آماری با وضعیت پایبندی کامل به اقدامات پیشگیری ثانویه نداشتند \\$(p>0.05)$.

\begin{landscape}
\begin{figure}[p]
\centering
\begin{tikzpicture}
\begin{axis}[
  xmode=log,
  log basis x=10,
  xmin=0.02,xmax=100,
  ymin=0.35,ymax=30.45,
  width=8.6cm,height=13.5cm,
  scale only axis,
  axis y line=none,
  axis x line*=bottom,
  ytick=\empty,
  xtick={0.02,0.05,0.1,0.2,0.5,1,2,5,10,20,50,100},
  xticklabels={0.02,0.05,0.1,0.2,0.5,1,2,5,10,20,50,100},
  tick label style={font=\scriptsize},
  xlabel={\rl{نسبت شانس (مقیاس لگاریتمی)} },
  xlabel style={font=\fontsize{10.0}{10.0}\selectfont,yshift=-1mm},
  xmajorgrids,
  major grid style={gray!25,line width=0.3pt},
  tick style={black,line width=0.5pt},
  clip=false
]
\draw[densely dashed,line width=0.9pt] (axis cs:1,0.5) -- (axis cs:1,29.5);
\node[anchor=south,fill=white,inner sep=1pt,font=\fontsize{9.0}{9.0}\selectfont]
  at (axis cs:1,29.55) {\rl{خط عدم اثر}};
\draw[gray!55,line width=0.45pt] ([xshift=-8.0cm]axis cs:0.02,26.5) -- ([xshift=5.0cm]axis cs:100,26.5);
\draw[gray!55,line width=0.45pt] ([xshift=-8.0cm]axis cs:0.02,18.5) -- ([xshift=5.0cm]axis cs:100,18.5);
\draw[gray!55,line width=0.45pt] ([xshift=-8.0cm]axis cs:0.02,13.5) -- ([xshift=5.0cm]axis cs:100,13.5);
\draw[gray!55,line width=0.45pt] ([xshift=-8.0cm]axis cs:0.02,6.5) -- ([xshift=5.0cm]axis cs:100,6.5);
\node[anchor=south east,align=center,text width=2.4cm,font=\bfseries\fontsize{11.0}{11.0}\selectfont]
  at ([xshift=-6.1cm]axis cs:0.02,30.15) {\rl{متغیر وابسته}};
\node[anchor=south east,align=center,text width=5.6cm,font=\bfseries\fontsize{11.0}{11.0}\selectfont]
  at ([xshift=-2mm]axis cs:0.02,30.15) {\rl{متغیر مستقل}};
\node[anchor=south,font=\bfseries\fontsize{11.0}{11.0}\selectfont]
  at (axis cs:1,30.15) {\rl{نسبت شانس و فاصله اطمینان ۹۵ درصد}};
\node[anchor=south west,align=center,text width=3.7cm,font=\bfseries\fontsize{10.0}{10.0}\selectfont]
  at ([xshift=2mm]axis cs:100,30.15) {\rl{نسبت شانس؛ فاصله اطمینان ۹۵ درصد}};
\node[anchor=south west,font=\bfseries\fontsize{13.0}{13.0}\selectfont]
  at ([xshift=4.1cm]axis cs:100,30.15) {\lr{p-value}};
\addplot[
  only marks,mark=square*,mark size=2.1pt,
  draw=black,fill=white,line width=0.7pt,
  error bars/.cd,x dir=both,x explicit,
  error bar style={line width=0.7pt},
  error mark options={rotate=90,mark size=2.0pt,line width=0.7pt}
] coordinates {
(0.98,29) += (0.070,0) -= (0.070,0)
(0.90,27) += (2.540,0) -= (0.670,0)
(0.96,26) += (0.090,0) -= (0.080,0)
(0.20,25) += (1.070,0) -= (0.180,0)
(2.84,24) += (19.050,0) -= (2.460,0)
(0.87,22) += (0.910,0) -= (0.440,0)
(0.30,21) += (1.750,0) -= (0.270,0)
(1.19,20) += (0.280,0) -= (0.200,0)
(0.40,19) += (1.500,0) -= (0.330,0)
(2.658,18) += (8.525,0) -= (2.026,0)
(1.062,17) += (1.731,0) -= (0.658,0)
(0.830,15) += (1.244,0) -= (0.498,0)
(0.585,14) += (0.790,0) -= (0.336,0)
(0.986,13) += (0.055,0) -= (0.051,0)
(2.109,10) += (4.077,0) -= (1.390,0)
(0.646,9) += (1.049,0) -= (0.400,0)
(0.878,8) += (1.777,0) -= (0.587,0)
(0.368,6) += (1.056,0) -= (0.273,0)
(7.824,5) += (61.899,0) -= (6.946,0)
(1.403,3) += (3.587,0) -= (1.009,0)
(0.941,2) += (0.160,0) -= (0.136,0)
(0.386,1) += (1.280,0) -= (0.297,0)
};
\addplot[
  only marks,mark=square*,mark size=2.4pt,
  draw=black,fill=black,line width=0.7pt,
  error bars/.cd,x dir=both,x explicit,
  error bar style={line width=0.9pt},
  error mark options={rotate=90,mark size=2.1pt,line width=0.9pt}
] coordinates {
(0.28,28) += (0.640,0) -= (0.210,0)
(0.19,23) += (0.670,0) -= (0.160,0)
(4.546,16) += (7.579,0) -= (2.842,0)
(0.232,12) += (0.415,0) -= (0.149,0)
(3.249,11) += (7.129,0) -= (2.232,0)
(0.797,7) += (0.115,0) -= (0.101,0)
(6.574,4) += (15.718,0) -= (4.635,0)
};
\ForestGroup{28}{\rl{قطع سیگار}}
\ForestGroup{22.5}{\rl{فعالیت بدنی}}
\ForestGroup{16}{\rl{رژیم غذایی}}
\ForestGroup{10}{\rl{سلامت روان‌شناختی}}
\ForestGroup{3.5}{\rl{پایبندی کلی}}
\ForestRow{29}{سن}{0.98}{0.91--1.05}{0.587}
\ForestRowSig{28}{\rl{اشتغال (شاغل)}}{0.28}{0.07--0.92}{0.037}
\ForestRow{27}{\rl{سواد سلامت (ناکافی)}}{0.90}{0.23--3.44}{0.886}
\ForestRow{26}{سن}{0.96}{0.88--1.05}{0.451}
\ForestRow{25}{\rl{تأهل (متأهل)}}{0.20}{0.02--1.27}{0.089}
\ForestRow{24}{\rl{تحصیلات (دانشگاهی)}}{2.84}{0.38--21.89}{0.298}
\ForestRowSig{23}{\rl{نوع درمان (دارویی در برابر آنژیوپلاستی)}}{0.19}{0.03--0.86}{0.031}
\ForestRow{22}{\rl{بیماری‌های همراه}}{0.87}{0.43--1.78}{0.707}
\ForestRow{21}{\rl{سواد سلامت (ناکافی)}}{0.30}{0.03--2.05}{0.220}
\ForestRow{20}{\rl{درک بیماری}}{1.19}{0.99--1.47}{0.062}
\ForestRow{19}{\rl{پایبندی دارویی (پایین)}}{0.40}{0.07--1.90}{0.257}
\ForestRow{18}{\rl{تحصیلات (دانشگاهی)}}{2.658}{0.632--11.183}{0.182}
\ForestRow{17}{\rl{سکونت (شهری)}}{1.062}{0.404--2.793}{0.903}
\ForestRowSig{16}{\rl{وضعیت اقتصادی (ضعیف)}}{4.546}{1.704--12.125}{0.002}
\ForestRow{15}{\rl{سواد سلامت (ناکافی)}}{0.830}{0.332--2.074}{0.690}
\ForestRow{14}{\rl{پایبندی دارویی (پایین)}}{0.585}{0.249--1.375}{0.219}
\ForestRow{13}{سن}{0.986}{0.935--1.041}{0.609}
\ForestRowSig{12}{\rl{جنسیت (زن)}}{0.232}{0.083--0.647}{0.005}
\ForestRowSig{11}{\rl{تأهل (متأهل)}}{3.249}{1.017--10.378}{0.047}
\ForestRow{10}{\rl{تحصیلات (دانشگاهی)}}{2.109}{0.719--6.186}{0.174}
\ForestRow{9}{\rl{اشتغال (شاغل)}}{0.646}{0.246--1.695}{0.374}
\ForestRow{8}{\rl{سواد سلامت (ناکافی)}}{0.878}{0.291--2.655}{0.818}
\ForestRowSig{7}{\rl{درک بیماری}}{0.797}{0.696--0.912}{0.001}
\ForestRow{6}{\rl{جنسیت (زن)}}{0.368}{0.095--1.424}{0.148}
\ForestRow{5}{\rl{تأهل (متأهل)}}{7.824}{0.878--69.723}{0.065}
\ForestRowSig{4}{\rl{تحصیلات (دانشگاهی)}}{6.574}{1.939--22.292}{0.003}
\ForestRow{3}{\rl{سواد سلامت (ناکافی)}}{1.403}{0.394--4.990}{0.601}
\ForestRow{2}{\rl{درک بیماری}}{0.941}{0.805--1.101}{0.450}
\ForestRow{1}{\rl{پایبندی دارویی (پایین)}}{0.386}{0.089--1.666}{0.202}
\end{axis}
\end{tikzpicture}
\caption{فارست پلات نتایج تحلیل رگرسیون لجستیک چندمتغیره}
\label{fig:multivariable-forest}
\end{figure}
\end{landscape}